\documentclass[UTF8]{ctexart}
\usepackage{amsmath, amssymb, geometry}
\usepackage{xcolor}
\usepackage{enumitem}
\usepackage{tcolorbox}
\usepackage{cases}
\usepackage{fancyhdr}
\geometry{a4paper, margin=1in}

% 定义红色环境
\newenvironment{redenv}{\color{red}}{}

% 定义详细解析环境
\newenvironment{detailedanalysis}{%
    \color{red}
    \begin{enumerate}[label=\textbf{第\arabic*步：}, leftmargin=*, align=left, itemsep=1.5ex]
}{%
    \end{enumerate}
}

% 定义概念框
\newenvironment{conceptbox}{%
    \begin{tcolorbox}[colback=red!5, colframe=red!50!black, title=概念解析, fonttitle=\bfseries]
}{%
    \end{tcolorbox}
}

% 定义公式推导环境
\newenvironment{formuladerivation}{%
    \begin{enumerate}[label={\textbf{公式\arabic*：}}, leftmargin=*, align=left]
}{%
    \end{enumerate}
}

% 定义题目和解析的分隔线
\newcommand{\problemrule}{\noindent\rule{\textwidth}{1.5pt}\vspace{-0.8em}}
\newcommand{\analysisrule}{\noindent\rule{\textwidth}{0.8pt}\vspace{-0.8em}}

% 宏定义：方便排版选择题选项
\newcommand{\choicesFour}[4]{
    \begin{enumerate}[label=\Alph*., itemsep=0pt, parsep=0pt, topsep=5pt, labelsep=0.5em]
        \begin{minipage}{0.22\linewidth} \item #1 \end{minipage}
        \begin{minipage}{0.22\linewidth} \item #2 \end{minipage}
        \begin{minipage}{0.22\linewidth} \item #3 \end{minipage}
        \begin{minipage}{0.22\linewidth} \item #4 \end{minipage}
    \end{enumerate}
}
\newcommand{\choicesTwo}[4]{
    \begin{enumerate}[label=\Alph*., itemsep=0pt, parsep=0pt, topsep=5pt, labelsep=0.5em]
        \begin{minipage}{0.45\linewidth} \item #1 \end{minipage}
        \begin{minipage}{0.45\linewidth} \item #2 \end{minipage}
        \begin{minipage}{0.45\linewidth} \item #3 \end{minipage}
        \begin{minipage}{0.45\linewidth} \item #4 \end{minipage}
    \end{enumerate}
}
\newcommand{\choices}[4]{
    \begin{enumerate}[label=\Alph*., itemsep=0pt, parsep=0pt, topsep=5pt, labelsep=0.5em]
        \item #1
        \item #2
        \item #3
        \item #4
    \end{enumerate}
}

\newcommand{\blank}[1]{\underline{\hspace{#1}}}

\begin{document}

\section*{第五章 中值定理与导数（提高）}

% ========== 第1题 ==========
\problemrule
\section*{【题目1】渐近线判定}
曲线 $ y=\frac{1}{f(x)} $ 有水平渐近线的充分条件是 \blank{1cm}。
\choicesTwo
{$ \lim_{x\to\infty}f(x)=0 $}
{$ \lim_{x\to\infty}f(x)=\infty $}
{$ \lim_{x\to0}f(x)=0 $}
{$ \lim_{x\to0}f(x)=\infty $}

\begin{redenv}
\analysisrule
\subsection*{逐步解析}

\begin{detailedanalysis}
\item \textbf{详细求解过程}
水平渐近线定义：若 $\lim_{x \to \infty} y = C$（常数），则 $y=C$ 是水平渐近线。
对应本题 $y = \frac{1}{f(x)}$，即要求 $\lim_{x \to \infty} \frac{1}{f(x)} = C$。
\begin{enumerate}[label=\Alph*.]
\item 若 $\lim_{x \to \infty} f(x) = 0$，则 $\lim_{x \to \infty} \frac{1}{f(x)} = \infty$，无水平渐近线。
\item 若 $\lim_{x \to \infty} f(x) = \infty$，则 $\lim_{x \to \infty} \frac{1}{f(x)} = 0$，有水平渐近线 $y=0$。
\item C, D 选项考查 $x \to 0$，这与水平渐近线（$x \to \infty$）无关。
\end{enumerate}

\item \textbf{最终答案}
\textbf{答案}：B
\end{detailedanalysis}
\end{redenv}

% ========== 第2题 ==========
\problemrule
\section*{【题目2】渐近线计数}
曲线 $ y = \frac{x}{1 - x^{2}} $ 的渐近线有 \blank{1cm}。
\choicesFour{1 条}{2 条}{3 条}{4 条}

\begin{redenv}
\analysisrule
\subsection*{逐步解析}

\begin{detailedanalysis}
\item \textbf{详细求解过程}
\begin{enumerate}[label=(\arabic*)]
\item \textbf{垂直渐近线}：
分母为零的点：$1 - x^2 = 0 \implies x = 1, x = -1$。
$\lim_{x \to 1} \frac{x}{1-x^2} = \infty$，$\lim_{x \to -1} \frac{x}{1-x^2} = \infty$。
所以有两条垂直渐近线：$x=1$ 和 $x=-1$。
\item \textbf{水平渐近线}：
$\lim_{x \to \infty} \frac{x}{1-x^2} = 0$。
所以有一条水平渐近线：$y=0$。
\item \textbf{斜渐近线}：
有水平渐近线通常不再有斜渐近线（双侧趋向时）。
\end{enumerate}
共 3 条。

\item \textbf{最终答案}
\textbf{答案}：C
\end{detailedanalysis}
\end{redenv}

% ========== 第3题 ==========
\problemrule
\section*{【题目3】拐点性质}
点 $(1,2)$ 是 $ f(x)=(x-a)^{3}+b $ 对应图形的拐点，则 \blank{1cm}。
\choicesFour{$ a=0,b=1 $}{$ a=2,b=3 $}{$ a=1,b=2 $}{$ a=-1,b=-6 $}

\begin{redenv}
\analysisrule
\subsection*{逐步解析}

\begin{detailedanalysis}
\item \textbf{详细求解过程}
\begin{enumerate}[label=(\arabic*)]
\item \textbf{求二阶导数}：
$f'(x) = 3(x-a)^2$
$f''(x) = 6(x-a)$
\item \textbf{拐点条件}：
拐点处 $f''(x)=0$（或不存在）。令 $6(x-a)=0 \implies x=a$。
题目已知拐点横坐标为 $1$，所以 $a=1$。
\item \textbf{点在曲线上}：
将 $(1,2)$ 代入方程：$f(1) = (1-a)^3 + b = 2$。
已知 $a=1$，则 $(1-1)^3 + b = 2 \implies b = 2$。
\end{enumerate}

\item \textbf{最终答案}
\textbf{答案}：C
\end{detailedanalysis}
\end{redenv}

% ========== 第4题 ==========
\problemrule
\section*{【题目4】拐点坐标}
函数 $ y = |\ln x| $ 对应图形的拐点是 \blank{1cm}。
\choicesFour{$ (1,0) $}{$ (e,1) $}{$ (2,\ln 2) $}{不存在}

\begin{redenv}
\analysisrule
\subsection*{逐步解析}

\begin{detailedanalysis}
\item \textbf{详细求解过程}
$$ y = \begin{cases} -\ln x, & 0 < x < 1 \\ \ln x, & x \ge 1 \end{cases} $$
\begin{enumerate}[label=(\arabic*)]
\item \textbf{求导数}：
$y' = \begin{cases} -1/x, & 0<x<1 \\ 1/x, & x>1 \end{cases}$
$y'' = \begin{cases} 1/x^2, & 0<x<1 \\ -1/x^2, & x>1 \end{cases}$
\item \textbf{凹凸性分析}：
$x \in (0, 1)$，$y'' > 0$，凹。
$x \in (1, \infty)$，$y'' < 0$，凸。
\item \textbf{拐点}：
在 $x=1$ 处连续，且凹凸性改变，所以 $(1, 0)$ 是拐点。
(注意：虽然 $x=1$ 处不可导，但满足连续曲线且凹凸方向改变的定义，通常视为拐点，或角点兼拐点)。
严格来说，拐点定义要求该点处有切线或连续。本题选项中有 $(1,0)$，选之。
\end{enumerate}

\item \textbf{最终答案}
\textbf{答案}：A
\end{detailedanalysis}
\end{redenv}

% ========== 第5题 ==========
\problemrule
\section*{【题目5】Rolle 定理推论}
设函数 $ f(x)=(x-1)(x-2)(x-3) $ ，则方程 $ f'(x)=0 $ 有 \blank{1cm}。
\choicesFour{一个实根}{二个实根}{三个实根}{无实根}

\begin{redenv}
\analysisrule
\subsection*{逐步解析}

\begin{detailedanalysis}
\item \textbf{详细求解过程}
\begin{enumerate}[label=(\arabic*)]
\item \textbf{零点分析}：$f(x)$ 有三个零点：$1, 2, 3$。
\item \textbf{Rolle 定理应用}：
在 $[1, 2]$ 上应用 Rolle 定理，存在 $\xi_1 \in (1, 2)$ 使得 $f'(\xi_1) = 0$。
在 $[2, 3]$ 上应用 Rolle 定理，存在 $\xi_2 \in (2, 3)$ 使得 $f'(\xi_2) = 0$。
\item \textbf{根的个数}：
因为 $f(x)$ 是三次多项式，导数 $f'(x)$ 是二次多项式，最多有两个根。
已找到两个不同实根，故恰有两个实根。
\end{enumerate}

\item \textbf{最终答案}
\textbf{答案}：B
\end{detailedanalysis}
\end{redenv}

% ========== 第6题 ==========
\problemrule
\section*{【题目6】极值与拐点判定}
设函数 $ f(x) = x^{3} + ax^{2} + bx + c $ , 且 $ f(0) = f'(0) = 0 $ , 则下列结论不正确的是 \blank{1cm}。
\choicesTwo
{$ b = c = 0 $}
{当 $a > 0$ 时， $ f(0) $ 为极小值}
{当 $a < 0$ 时， $ f(0) $ 为极大值}
{当 $ a \neq 0 $ 时， $ (0, f(0)) $ 为拐点}

\begin{redenv}
\analysisrule
\subsection*{逐步解析}

\begin{detailedanalysis}
\item \textbf{详细求解过程}
\begin{enumerate}[label=(\arabic*)]
\item \textbf{利用条件确定参数}：
$f(0) = c = 0 \implies c=0$。
$f'(x) = 3x^2 + 2ax + b$，
$f'(0) = b = 0 \implies b=0$。
所以 $f(x) = x^3 + ax^2$。选项 A 正确。
\item \textbf{极值判定}：
$f'(x) = 3x^2 + 2ax = x(3x + 2a)$。驻点为 $0, -2a/3$。
$f''(x) = 6x + 2a$。
$f''(0) = 2a$。
- 若 $a > 0$，则 $f''(0) > 0$，为极小值。B 正确。
- 若 $a < 0$，则 $f''(0) < 0$，为极大值。C 正确。
\item \textbf{拐点判定}：
拐点处 $f''(x) = 0 \implies 6x + 2a = 0 \implies x = -a/3$。
当 $a \neq 0$ 时，$x = -a/3 \neq 0$。
所以在 $x=0$ 处 $f''(0) = 2a \neq 0$，不是拐点。D 错误。
\end{enumerate}

\item \textbf{最终答案}
\textbf{答案}：D
\end{detailedanalysis}
\end{redenv}

% ========== 第7题 ==========
\problemrule
\section*{【题目7】单调区间求解}
函数 $ y=\frac{x}{\ln x} $ 的单调增加的区间是 \blank{2cm}。

\begin{redenv}
\analysisrule
\subsection*{逐步解析}

\begin{detailedanalysis}
\item \textbf{详细求解过程}
\begin{enumerate}[label=(\arabic*)]
\item \textbf{定义域}：$x > 0$ 且 $\ln x \neq 0 \implies x \neq 1$。
定义域为 $(0, 1) \cup (1, +\infty)$。
\item \textbf{求导}：
$$ y' = \frac{1 \cdot \ln x - x \cdot \frac{1}{x}}{(\ln x)^2} = \frac{\ln x - 1}{(\ln x)^2} $$
\item \textbf{判定符号}：
令 $y' > 0 \implies \ln x - 1 > 0 \implies \ln x > 1 \implies x > e$。
所以单调增区间为 $(e, +\infty)$。
\end{enumerate}

\item \textbf{最终答案}
\textbf{答案}：$ (e, +\infty) $
\end{detailedanalysis}
\end{redenv}

% ========== 第8题 ==========
\problemrule
\section*{【题目8】驻点求解}
函数 $ y=(x-2)^{3} $ 的驻点是 \blank{2cm}。

\begin{redenv}
\analysisrule
\subsection*{逐步解析}

\begin{detailedanalysis}
\item \textbf{详细求解过程}
驻点即 $y' = 0$ 的点。
$y' = 3(x-2)^2$。
令 $y' = 0 \implies x - 2 = 0 \implies x = 2$。

\item \textbf{最终答案}
\textbf{答案}：$ x=2 $
\end{detailedanalysis}
\end{redenv}

% ========== 第9题 ==========
\problemrule
\section*{【题目9】凹凸性区间}
函数 $ y=-x^{3}+x^{2} $ 的凸区间是 \blank{2cm}。

\begin{redenv}
\analysisrule
\subsection*{逐步解析}

\begin{detailedanalysis}
\item \textbf{问题本质分析}
凸区间通常指 $y'' < 0$ 的区间（注：不同教材对凹凸定义可能相反，通常"凸"指向上凸，即 $\cap$ 形，对应二阶导数小于0）。
\item \textbf{详细求解过程}
\begin{enumerate}[label=(\arabic*)]
\item \textbf{求导}：
$y' = -3x^2 + 2x$
$y'' = -6x + 2$
\item \textbf{解不等式}：
令 $y'' < 0 \implies -6x + 2 < 0 \implies 2 < 6x \implies x > \frac{1}{3}$。
注：若教材定义下凸为凸，则解 $y'' > 0$。根据常规数学分析（同济版），图形向上凸（凸弧）对应 $y'' < 0$。
所以区间为 $(\frac{1}{3}, +\infty)$。
\end{enumerate}

\item \textbf{最终答案}
\textbf{答案}：$ (\frac{1}{3}, +\infty) $
\end{detailedanalysis}
\end{redenv}

% ========== 第10题 ==========
\problemrule
\section*{【题目10】垂直渐近线}
曲线 $ y = \frac{x^{2}-2x+2}{x-1} $ 的垂直渐近线的方程是 \blank{2cm}。

\begin{redenv}
\analysisrule
\subsection*{逐步解析}

\begin{detailedanalysis}
\item \textbf{详细求解过程}
分母为零点 $x=1$。
分子在 $x=1$ 时：$1 - 2 + 2 = 1 \neq 0$。
所以 $\lim_{x \to 1} y = \infty$。
故 $x=1$ 是垂直渐近线。

\item \textbf{最终答案}
\textbf{答案}：$ x=1 $
\end{detailedanalysis}
\end{redenv}

% ========== 第11题 ==========
\problemrule
\section*{【题目11】拐点计算}
求曲线 $ y = x^{3} - 3x^{2} + 3 $ 的拐点 \blank{2cm}。

\begin{redenv}
\analysisrule
\subsection*{逐步解析}

\begin{detailedanalysis}
\item \textbf{详细求解过程}
\begin{enumerate}[label=(\arabic*)]
\item \textbf{求导}：
$y' = 3x^2 - 6x$
$y'' = 6x - 6$
\item \textbf{找零点}：
令 $y'' = 0 \implies x = 1$。
\item \textbf{验证变号}：
$x < 1$ 时 $y'' < 0$；$x > 1$ 时 $y'' > 0$。变号，故为拐点。
\item \textbf{求坐标}：
$y(1) = 1 - 3 + 3 = 1$。
拐点为 $(1, 1)$。
\end{enumerate}

\item \textbf{最终答案}
\textbf{答案}：$ (1,1) $
\end{detailedanalysis}
\end{redenv}

% ========== 第12题 ==========
\problemrule
\section*{【题目12】Lagrange 中值定理辅助函数}
若 $ f(x) $ 在 $ [0,1] $ 上连续，在 $ (0,1) $ 内可导，且 $ f(0)=f(1)=0, f(1)=1 $ (题目此处矛盾，应为 $f(0)=0, f(1)=1$)。
证明在 $ (0,1) $ 内至少有一点 $ \xi $ 使 $ f'(\xi)=1 $ 。

\begin{redenv}
\analysisrule
\subsection*{逐步解析}

\begin{detailedanalysis}
\item \textbf{题目修正}
题目既然要证 $f'(\xi)=1$，且 $f(0)$ 和 $f(1)$ 给定。结合 $f(0)=0, f(1)=1$，直接用拉格朗日中值定理即可。原题写 $f(0)=f(1)=0, f(1)=1$ 显然笔误。应理解为 $f(0)=0, f(1)=1$。
\item \textbf{证明过程}
在区间 $[0, 1]$ 上应用拉格朗日中值定理：
存在 $\xi \in (0, 1)$，使得
$$ f'(\xi) = \frac{f(1) - f(0)}{1 - 0} = \frac{1 - 0}{1} = 1 $$
证毕。
\end{detailedanalysis}
\end{redenv}

% ========== 第13题 ==========
\problemrule
\section*{【题目13】最值问题}
求斜边长为定长 $L$ 的直角三角形的最大面积。

\begin{redenv}
\analysisrule
\subsection*{逐步解析}

\begin{detailedanalysis}
\item \textbf{详细求解过程}
设一直角边为 $x$，则另一直角边为 $\sqrt{L^2 - x^2}$ ($0 < x < L$)。
\begin{enumerate}[label=(\arabic*)]
\item \textbf{建立面积函数}：
$S(x) = \frac{1}{2} x \sqrt{L^2 - x^2}$。
为方便计算，求 $S^2(x)$ 的最大值。
$f(x) = S^2(x) = \frac{1}{4} x^2 (L^2 - x^2) = \frac{1}{4} (L^2 x^2 - x^4)$。
\item \textbf{求导}：
$f'(x) = \frac{1}{4} (2L^2 x - 4x^3) = \frac{x}{2} (L^2 - 2x^2)$。
令 $f'(x) = 0$，得 $2x^2 = L^2 \implies x = \frac{L}{\sqrt{2}}$。
\item \textbf{求最大面积}：
此时另一边也是 $\frac{L}{\sqrt{2}}$，即等腰直角三角形。
面积 $S = \frac{1}{2} \cdot \frac{L}{\sqrt{2}} \cdot \frac{L}{\sqrt{2}} = \frac{L^2}{4}$。
\end{enumerate}

\item \textbf{最终答案}
\textbf{答案}：$ \frac{L^2}{4} $
\end{detailedanalysis}
\end{redenv}

% ========== 第14题 ==========
\problemrule
\section*{【题目14】函数性质全面讨论}
讨论 $ y = xe^{-x} $ 的增减性、凹凸性、极值、拐点。

\begin{redenv}
\analysisrule
\subsection*{逐步解析}

\begin{detailedanalysis}
\item \textbf{详细求解过程}
\begin{enumerate}[label=(\arabic*)]
\item \textbf{一阶导数（单调性与极值）}：
$y' = e^{-x} - xe^{-x} = e^{-x}(1-x)$。
令 $y'=0 \implies x=1$。
- $x < 1$: $y' > 0$，单调递增。
- $x > 1$: $y' < 0$，单调递减。
极值：$x=1$ 处有极大值 $y(1) = e^{-1} = \frac{1}{e}$。
\item \textbf{二阶导数（凹凸性与拐点）}：
$y'' = (e^{-x}(1-x))' = -e^{-x}(1-x) + e^{-x}(-1) = e^{-x}(x-1-1) = e^{-x}(x-2)$。
令 $y''=0 \implies x=2$。
- $x < 2$: $y'' < 0$，凸（向上凸）。
- $x > 2$: $y'' > 0$，凹（向上凹）。
拐点：$(2, 2e^{-2})$。
\end{enumerate}

\item \textbf{最终答案}
\textbf{增区间}：$(-\infty, 1)$；\textbf{减区间}：$(1, +\infty)$。
\textbf{极大值}：$1/e$；无极小值。
\textbf{凸区间}：$(-\infty, 2)$；\textbf{凹区间}：$(2, +\infty)$。
\textbf{拐点}：$(2, 2e^{-2})$。
\end{detailedanalysis}
\end{redenv}

% ========== 第15题 ==========
\problemrule
\section*{【题目15】经济应用}
已知某产品的需求函数 $ Q = 28 - 2p $ 和总成本函数 $ C(Q) = 30 + 2Q $ , 其中 $Q$ 为销售量，$p$ 为价格，求当 $p$ 为多少可获得最大利润？

\begin{redenv}
\analysisrule
\subsection*{逐步解析}

\begin{detailedanalysis}
\item \textbf{详细求解过程}
\begin{enumerate}[label=(\arabic*)]
\item \textbf{建立利润函数}：
利润 $L = \text{收入} - \text{成本} = R - C$。
收入 $R = p \cdot Q = p(28 - 2p) = 28p - 2p^2$。
成本 $C = 30 + 2Q = 30 + 2(28 - 2p) = 30 + 56 - 4p = 86 - 4p$。
利润 $L(p) = (28p - 2p^2) - (86 - 4p) = -2p^2 + 32p - 86$。
\item \textbf{求极值}：
求导：$L'(p) = -4p + 32$。
令 $L'(p) = 0 \implies 4p = 32 \implies p = 8$。
二阶导 $L''(p) = -4 < 0$，确实是最大值点。
\end{enumerate}

\item \textbf{最终答案}
\textbf{答案}：$ p=8 $
\end{detailedanalysis}
\end{redenv}

% ========== 第16题 (原题15+1) ==========
\problemrule
\section*{【题目16】不等式证明}
证明不等式 $ e^{x} < (1 + x)^{(1 + x)} $ ($x > 0$)。

\begin{redenv}
\analysisrule
\subsection*{逐步解析}

\begin{detailedanalysis}
\item \textbf{详细求解过程}
两边取对数（因均为正）：
需证 $x < (1+x) \ln(1+x)$。
即证 $f(x) = (1+x) \ln(1+x) - x > 0$ 当 $x > 0$ 时。
\begin{enumerate}[label=(\arabic*)]
\item \textbf{求导}：
$f'(x) = 1 \cdot \ln(1+x) + (1+x) \cdot \frac{1}{1+x} - 1 = \ln(1+x) + 1 - 1 = \ln(1+x)$。
\item \textbf{单调性}：
当 $x > 0$ 时，$\ln(1+x) > \ln 1 = 0$。
所以 $f(x)$ 在 $(0, +\infty)$ 上单调递增。
\item \textbf{比较初值}：
$f(0) = 1 \cdot 0 - 0 = 0$。
所以当 $x > 0$ 时，$f(x) > f(0) = 0$。
即原不等式成立。
\end{enumerate}
\end{detailedanalysis}
\end{redenv}

\end{document}
